\documentclass[11pt,a4paper]{article}

% ──────────────────────────────────────────────────────────────────────────
% EasyNet Invocation Security Model — Reference Specification
% Status: Industrial-Grade Target Design (design-first, implementation-agnostic)
% Author: Silan Hu
% Build:   xelatex invocation_security_v1.tex   (CJK + Unicode required)
%
% Sibling document to easynet_ontology.tex in the same directory.
% Shares vocabulary, style, and the [axiom]/[derived]/[deferred] discipline.
%
% This document is design-first. The body (§§1--13) describes the target
% security system as if it were being defined from scratch for a reference
% implementation. It does not reference the current code at all. Current
% implementation status is covered exclusively in the appendices.
% ──────────────────────────────────────────────────────────────────────────

\usepackage[margin=1in]{geometry}
\usepackage{amsmath,amssymb,amsthm}
\usepackage{xcolor}
\usepackage[colorlinks=true,linkcolor=blue!60!black,citecolor=blue!60!black,urlcolor=blue!60!black]{hyperref}
\usepackage{booktabs}
\usepackage{longtable}
\usepackage{array}
\usepackage{tabularx}
\usepackage{xltabular}
\newcolumntype{L}[1]{>{\raggedright\arraybackslash}p{#1}}
\usepackage{listings}
\usepackage{enumitem}
\usepackage{tcolorbox}
\usepackage{titlesec}
\usepackage{fancyhdr}
\usepackage{parskip}
\usepackage{xspace}

% ── Unicode + CJK support (xelatex) ───────────────────────────────────────
\usepackage{fontspec}
\usepackage{xeCJK}
\setCJKmainfont{PingFang SC}[
    BoldFont={PingFang SC Semibold},
    ItalicFont={PingFang SC Light}
]

% ── Page style ────────────────────────────────────────────────────────────
\pagestyle{fancy}
\fancyhf{}
\fancyhead[L]{\small EasyNet Invocation Security Model}
\fancyhead[R]{\small Reference Specification}
\fancyfoot[C]{\thepage}
\renewcommand{\headrulewidth}{0.4pt}

% ── Section styling ───────────────────────────────────────────────────────
\titleformat{\section}{\Large\bfseries\color{black!85}}{\thesection}{0.7em}{}
\titleformat{\subsection}{\large\bfseries\color{black!75}}{\thesubsection}{0.6em}{}
\titleformat{\subsubsection}{\normalsize\bfseries\color{black!70}}{\thesubsubsection}{0.5em}{}

% ── Tag colors ────────────────────────────────────────────────────────────
\definecolor{axiomcolor}{RGB}{180,30,30}
\definecolor{derivedcolor}{RGB}{30,90,180}
\definecolor{deferredcolor}{RGB}{120,90,30}
\definecolor{defectcolor}{RGB}{150,20,80}
\definecolor{codebg}{RGB}{245,245,248}
\definecolor{codeborder}{RGB}{210,210,220}

% ── Tag macros ────────────────────────────────────────────────────────────
\newcommand{\axiom}{\textcolor{axiomcolor}{\textbf{[axiom]}}\xspace}
\newcommand{\derived}{\textcolor{derivedcolor}{\textbf{[derived]}}\xspace}
\newcommand{\deferred}{\textcolor{deferredcolor}{\textbf{[deferred]}}\xspace}
\newcommand{\defect}{\textcolor{defectcolor}{\textbf{[non-conformance]}}\xspace}

% ── Code listing style ────────────────────────────────────────────────────
\lstdefinestyle{eal}{
    backgroundcolor=\color{codebg},
    basicstyle=\ttfamily\small,
    frame=single,
    rulecolor=\color{codeborder},
    breaklines=true,
    showstringspaces=false,
    keywordstyle=\color{blue!70!black}\bfseries,
    commentstyle=\color{green!40!black}\itshape,
    stringstyle=\color{red!60!black},
    morekeywords={fn,pub,let,match,if,else,return,struct,impl,use,crate,self,message,enum},
    columns=fullflexible,
    keepspaces=true,
    aboveskip=0.6em,
    belowskip=0.6em,
}
\lstset{style=eal}

% ── Theorem-like environments ─────────────────────────────────────────────
\newtcolorbox{invariant}[1][]{
    colback=red!4!white,
    colframe=red!50!black,
    fonttitle=\bfseries,
    title={Invariant},
    sharp corners=south,
    boxrule=0.6pt,
    #1
}
\newtcolorbox{property}[1][]{
    colback=green!3!white,
    colframe=green!40!black,
    fonttitle=\bfseries,
    title={Security property},
    sharp corners=south,
    boxrule=0.6pt,
    #1
}
\newtcolorbox{notebox}[1][]{
    colback=blue!4!white,
    colframe=blue!50!black,
    fonttitle=\bfseries,
    title={Note},
    sharp corners=south,
    boxrule=0.6pt,
    #1
}
\newtcolorbox{conformancebox}[1][]{
    colback=orange!6!white,
    colframe=orange!70!black,
    fonttitle=\bfseries,
    title={Conformance gap},
    sharp corners=south,
    boxrule=0.6pt,
    #1
}

% ── TOC depth: section + subsection only. paragraph-level entries in the
%    conformance appendix must not pollute the table of contents.
\setcounter{tocdepth}{2}
\setcounter{secnumdepth}{2}

% ── Title ─────────────────────────────────────────────────────────────────
\title{
    \vspace{-2em}
    \textbf{\Large EasyNet Invocation Security Model}\\[0.3em]
    \large Reference Specification \textemdash{} Industrial-Grade Target
    Design
}
\author{
    Silan Hu \\ \small National University of Singapore
}
\date{Reference edition, \today}

% ──────────────────────────────────────────────────────────────────────────
\begin{document}
\maketitle

\begin{abstract}
\noindent
This document specifies the invocation security system for EasyNet at the
rigor of RFC-class reference specifications. The body is
\emph{design-first}: it defines the target system as if from a blank slate,
without reference to any current implementation. The central claim is that
an EasyNet invocation is a cryptographic \emph{statement} made by one or
more principals about an intended operation, and the runtime's sole
legitimate behavior is to evaluate that statement against a formally
defined binding semantics, fail-closed. The design is grounded in three
irreducible components: (i) per-layer cryptographic identity with
explicit binding rules (\S\ref{sec:binding}), (ii) an SSH-shaped
bootstrap separating single-use authorization tokens from durable
identity keypairs (\S\ref{sec:bootstrap}), and (iii) a mandatory
signing-and-verification pipeline (\S\S\ref{sec:sign}--\ref{sec:verify}).
It is aligned with and cites OpenSSH public-key authentication
(RFC~4252), TLS~1.3 (RFC~8446), FIDO2/WebAuthn (W3C), SPIFFE/SPIRE (CNCF),
NIST FIPS~186-5, FIPS~204, and the Dolev--Yao adversary model. Conformance
of any particular implementation to this specification is addressed
separately in the appendices.
\end{abstract}

\tableofcontents
\vspace{1em}
\hrule
\vspace{1em}

% ──────────────────────────────────────────────────────────────────────────
\section{Reading guide}
\label{sec:reading-guide}

This document uses the same three-tag discipline as
\texttt{easynet\_ontology.tex}, plus a fourth tag reserved for the
conformance appendix.

\begin{itemize}[leftmargin=2.4em,itemsep=0.3em]
    \item[\axiom] Either (i) a cited external standard (RFC, W3C,
    NIST, peer-reviewed paper) that this document adopts by reference,
    or (ii) a definition internal to this document that later
    derivations rest on. Axioms are load-bearing: changing one requires
    re-deriving every dependent statement.

    \item[\derived] A logical consequence of one or more axioms.
    Negotiable in review; the axioms are not.

    \item[\deferred] A known extension point, deliberately postponed to
    a later version of this document (\S\ref{sec:gaps}).

    \item[\defect] Reserved for the conformance appendix
    (\S\ref{app:conformance}). A \defect tag identifies a deviation of a
    specific implementation from this specification. It never appears
    in the body.
\end{itemize}

\begin{notebox}
The body of this document (\S\S\ref{sec:reading-guide}--\ref{sec:gaps}) is
design-first and self-contained: read in isolation, it is a complete
reference specification suitable for clean-room implementation. The
appendices (\S\ref{app:conformance} onward) analyze conformance of the
present EasyNet codebase against this specification, for operational
planning; they are not part of the design.
\end{notebox}

% ──────────────────────────────────────────────────────────────────────────
\section{Scope}
\label{sec:scope}

\axiom{} This specification governs the authentication, integrity, and
non-repudiation of \emph{invocations} within an EasyNet runtime. An
invocation is any control-plane or data-plane operation that names
(i) a tenant, (ii) a resource or capability to be exercised, and
(iii) zero or more principals claiming authorization for the operation.

\derived{} The following operations are in scope: node registration,
node heartbeat, ability invocation (synchronous, asynchronous,
streaming), capability lifecycle management, mission dispatch, and any
operation that returns an audit record.

\derived{} The following are out of scope for this document:

\begin{itemize}[leftmargin=2em,itemsep=0.2em]
    \item Confidentiality of invocation payloads. This is delegated to
    the underlying transport (TLS~1.3, RFC~8446~\cite{rfc8446}).
    \item Anonymity and unlinkability of callers.
    \item Network-level denial-of-service defense.
    \item Billing, metering, and quota enforcement.
    \item Application-layer authorization (what a principal is
    \emph{allowed} to invoke); this spec establishes the
    authenticated identity on which authorization may be built.
\end{itemize}

% ──────────────────────────────────────────────────────────────────────────
\section{Adversary model}
\label{sec:adversary}

\axiom{} The adversary is a strengthened Dolev--Yao
adversary~\cite{dolev1983security}, consistent with the threat models
of TLS~1.3 (RFC~8446~\cite{rfc8446} \S E), OpenSSH
(RFC~4252~\cite{rfc4252}), and SPIFFE~\cite{spiffe}:

\begin{enumerate}[leftmargin=2em,itemsep=0.2em]
    \item \textbf{Network control.} The adversary fully controls the
    network between any two parties: it may intercept, modify,
    replay, reorder, and inject arbitrary messages.
    \item \textbf{Selective corruption.} The adversary may corrupt any
    finite set of non-target parties and learn their long-term keys.
    \item \textbf{Cryptographic soundness.} The adversary cannot break
    Ed25519 signatures~\cite{bernstein2012high} or ECDSA-P256
    (FIPS~186-5~\cite{fips186-5}); cannot find SHA-256 preimages or
    collisions; cannot break the confidentiality, integrity, or
    authenticity of TLS~1.3.
    \item \textbf{Filesystem boundaries.} Private keys stored with
    correct permissions on parties not under the adversary's control
    remain private.
    \item \textbf{Public-key visibility.} The adversary may learn the
    full set of public keys the system has ever registered, including
    revoked ones.
\end{enumerate}

\derived{} Three consequences:

\begin{itemize}[leftmargin=2em,itemsep=0.2em]
    \item Channel-level mTLS alone is insufficient: a corrupted party
    with a valid channel certificate is indistinguishable from an
    honest one at the channel layer.
    \item Bearer credentials (any secret whose possession is the
    credential) are strictly weaker than proof-of-possession
    signatures.
    \item Server-side custody of a principal's private key reduces
    that principal's identity to a proxy of the custodian.
\end{itemize}

% ──────────────────────────────────────────────────────────────────────────
\section{Security goals}
\label{sec:goals}

\axiom{} This specification commits to the following six properties.
Each is formalized to the level of RFC~4252~\cite{rfc4252} \S7 and
WebAuthn~\cite{webauthn} \S13.

\begin{property}
\textbf{G1 (Authenticity).}
An invocation is accepted only if every identity it names is
authenticated by a signature verifiable against a public key registered
for that identity in state \texttt{Active}, under a private key the
identity owner alone possesses.
\end{property}

\begin{property}
\textbf{G2 (Integrity).}
Any modification, by a party other than the signer, of any field the
runtime uses to interpret the invocation
(tenant, resource, function, arguments, nonce, timestamp, principal
identifiers) invalidates the signature.
\end{property}

\begin{property}
\textbf{G3 (Non-repudiation).}
Every accepted invocation produces an audit record that independently
verifies against the registered public keys. A signer cannot credibly
deny having authorized an invocation whose signature verifies under its
key.
\end{property}

\begin{property}
\textbf{G4 (Replay-resistance).}
An invocation captured off the wire cannot be successfully re-submitted,
either within the nonce-reservation window (detected by nonce replay)
or outside it (detected by timestamp skew).
\end{property}

\begin{property}
\textbf{G5 (Forward security of audit).}
Compromise of a private key at time $t$ does not invalidate audit
records produced at $t' < t$ while the key was in state \texttt{Active}
and before revocation was applied.
\end{property}

\begin{property}
\textbf{G6 (Principle of least identity).}
An invocation asserts exactly the identities required to authorize the
statement it makes. One identity's credentials cannot stand in for
another's; delegation is an explicit, bounded, cryptographically
evidenced operation (binding rule B4, \S\ref{sec:gaps}).
\end{property}

\derived{} G1--G6 together define what ``authenticated'' means in this
specification. Confidentiality, anonymity, and unlinkability are
explicitly not listed; their absence from this list is intentional.

% ──────────────────────────────────────────────────────────────────────────
\section{The invocation as a cryptographic statement}
\label{sec:statement}

\subsection{Definition}

\axiom{} An invocation $I$ is a tuple
\[
I \;=\; (T,\,P,\,N,\,A,\,f,\,\mathsf{args},\,t_{\mathsf{iat}},\,
\mathsf{req\_id},\,\mathsf{nonce})
\]
where:
\begin{itemize}[leftmargin=2em,itemsep=0.15em]
    \item $T$ is the tenant under whose authority the invocation is
    made;
    \item $P$ is the set of principals claiming authorization (possibly
    empty);
    \item $N$ is the originating node (non-empty for any invocation
    that traverses an Axon runtime);
    \item $A$ is the ability or capability being exercised, and $f$
    its specific function;
    \item $\mathsf{args}$ is the argument payload;
    \item $t_{\mathsf{iat}}$ is the Unix-millisecond timestamp at
    which $I$ was issued;
    \item $\mathsf{req\_id}$ is a globally unique request identifier;
    \item $\mathsf{nonce}$ is at least 128 bits of cryptographic
    randomness.
\end{itemize}

\axiom{} The canonical byte encoding of $I$, denoted $C(I)$, is a
deterministic, injective serialization of the tuple under an
unambiguous delimiter scheme.\footnote{``Canonical'' in the sense of
RFC~8785~\cite{rfc8785} and the SSH signing-blob format in
RFC~4252~\cite{rfc4252} \S7: two parties computing $C(I)$ from the same
$I$ produce byte-identical results.}

\derived{} $C(I)$ includes $H(\mathsf{args})$ where $H$ is SHA-256, not
$\mathsf{args}$ itself; this keeps $C(I)$ bounded in size while
preserving integrity of arbitrarily large payloads.

\subsection{The authorization proposition}

\axiom{} Associated with $I$ is the \emph{authorization proposition}
$\Phi(I)$:

\begin{quote}
``Principal(s) $P$, acting under tenant $T$, transported by node $N$,
authorize the execution of function $f$ of ability $A$ with arguments
$\mathsf{args}$ at time $t_{\mathsf{iat}}$.''
\end{quote}

The runtime's sole task, at the invocation boundary, is to decide
whether $\Phi(I)$ is cryptographically established by the evidence
accompanying $I$.

% ──────────────────────────────────────────────────────────────────────────
\section{Identity model}
\label{sec:identity}

\subsection{Identity layers}

\axiom{} The specification recognizes three layers of cryptographic
identity:

\begin{itemize}[leftmargin=2em,itemsep=0.25em]
    \item \textbf{Node identity ($K_N$).} A keypair bound to a specific
    node (host, device, runtime instance). The private component is
    generated and retained on the node; the public component is
    registered with the runtime's identity service at pairing time.
    \item \textbf{Principal identity ($K_P$).} A keypair bound to a
    logical principal (human user, service account, agent). Two
    registration modes are recognized, distinguished by who holds the
    private component:
    \begin{itemize}
        \item \emph{Strong principal identity.} The private component is
        generated and held exclusively by the principal's own trust
        boundary (e.g.\ a user's browser via
        WebCrypto~\cite{webcrypto}; a developer's local machine).
        \item \emph{Weak principal identity.} The private component is
        generated and held by an intermediary acting on the principal's
        behalf (e.g.\ a backend service-account key signing
        ``on behalf of user $U$''). Weak identity is explicitly marked
        in the registration metadata and surfaced in the audit record.
    \end{itemize}
    \item \textbf{Trust root identity ($K_R$).} The keypair or
    certificate authority that endorses the runtime's public identity,
    permitting clients to verify the runtime itself. This is the
    runtime's counterpart to SSH's \texttt{known\_hosts}.
\end{itemize}

\axiom{} All identities use the same underlying signature suite (see
\S\ref{sec:sign}.3) so that evidence is uniform and pipelines are
algorithm-agile.

\subsection{Registration}

\axiom{} Every identity is associated with a \emph{key record} of the
form:
\begin{lstlisting}
KeyRecord {
    subject:    { kind: Node | Principal | Root, id: string, tenant: string },
    algorithm:  Ed25519 | ECDSA-P256 | ...,
    version:    u32,              // monotonically increasing
    public_key: bytes,
    state:      Active | Rotated | Revoked,
    mode:       Strong | Weak,    // principals only; nodes are always Strong
    created_at: int64,
    rotated_at: int64,
    expires_at: int64,            // 0 = no expiry
    fingerprint: bytes            // SHA-256 of public_key
}
\end{lstlisting}

\derived{} The \texttt{fingerprint} is redundant with
\texttt{public\_key} but provides a stable short identifier for audit
records.

\subsection{Binding semantics}
\label{sec:binding}

\axiom{} A signature $\sigma$ by private key $\mathit{sk}$ over
canonical bytes $C(I)$, verifying against public key
$K = \mathsf{pk}(\mathit{sk})$, authorizes the low-level proposition:

\begin{quote}
``The holder of the private key matching $K$ authorized the byte
sequence $C(I)$ at the time $C(I)$ claims.''
\end{quote}

Translating this into a statement about $\Phi(I)$ requires an explicit
binding rule.

\axiom{} This specification defines three binding rules. No other
binding is permitted for v1.

\begin{property}
\textbf{B1 (Node origin).}
A signature $\sigma_N$ under $K_N$ (a node's registered key) over
$C(I)$ authorizes the proposition ``\emph{invocation $I$ originates
from node $N$.}'' It asserts nothing about any principal.
\end{property}

\begin{property}
\textbf{B2 (Principal authorization).}
A signature $\sigma_P$ under $K_P$ (a principal's registered key) over
$C(I)$ authorizes the proposition ``\emph{principal $P$ personally
authorized invocation $I$.}'' Strong B2 obtains when $K_P$ is
registered as \texttt{Strong}; weak B2 obtains when \texttt{Weak}.
Both are B2; the distinction is evidentiary strength, not kind.
\end{property}

\begin{property}
\textbf{B3 (Conjunction).}
An invocation carrying both $\sigma_N$ and $\sigma_P$, each verifying
under its respective rule over the same $C(I)$, authorizes the
conjunction $\Phi(I)$ in full: $P$ personally authorized, $N$
transported.
\end{property}

\begin{invariant}
\textbf{(Binding non-substitutability.)} A signature established under
rule $B_i$ shall not count as evidence for rule $B_j$, $j \ne i$. In
particular, a node signature is not acceptable in lieu of a principal
signature for invocations that name a principal, and a principal
signature is not acceptable in lieu of a node signature for the node-
origin claim.
\end{invariant}

\derived{} When $\Phi(I)$ names $k$ principals (e.g.\ a multi-party
authorization), B2 is required $k$ times, once per principal, each
under the respective $K_P$. The general rule: the runtime accepts $I$
only when $\Phi(I)$ is covered by the conjunction of all required
bindings.

% ──────────────────────────────────────────────────────────────────────────
\section{Bootstrap: separating authorization from identity}
\label{sec:bootstrap}

\subsection{The problem bootstrap solves}

\axiom{} Identity registration is itself an operation that must be
authorized, yet the identity being registered does not yet exist at
authorization time. This chicken-and-egg problem is solved in
OpenSSH~\cite{rfc4252} by separating two roles: a short-lived
\emph{bootstrap authorization} (e.g.\ the password accepted by
\texttt{ssh-copy-id}), which authorizes exactly one transaction
---the placement of a public key into
\texttt{authorized\_keys}---and a \emph{durable identity} (the
keypair) that authenticates every subsequent session.

\subsection{Bootstrap tokens}

\axiom{} A \emph{bootstrap token} is a high-entropy, server-issued,
single-use credential. It satisfies:

\begin{itemize}[leftmargin=2em,itemsep=0.15em]
    \item \textbf{High entropy.} At least 256 bits of cryptographic
    randomness, rendered in a URL-safe encoding.
    \item \textbf{Single-use.} The server atomically consumes it on
    validation; a second presentation fails.
    \item \textbf{Short-lived.} Validity is bounded by an explicit
    expiry (recommended: $\le 15$ minutes from issuance).
    \item \textbf{Narrowly scoped.} Authorizes one transaction of one
    kind: the registration of a single keypair under a single named
    identity.
\end{itemize}

\axiom{} The bootstrap token authorizes the transaction ``register
public key $K$ as the durable identity of subject $S$ within tenant
$T$.'' That is its entire authority. It is not an identity and is not
accepted as evidence of any proposition about $I$ after the registration
transaction completes.

\subsection{Durable keypair registration}

\axiom{} Durable keypair registration proceeds as follows. Let the
prospective subject be $S$, the tenant $T$, and the bootstrap token
$\tau$:

\begin{enumerate}[leftmargin=2em,itemsep=0.2em]
    \item $S$ locally generates a fresh keypair
    $(\mathit{sk}_S, K_S)$ using the operating system's cryptographic
    RNG (\S\ref{sec:keys}.1).
    \item $S$ computes a \emph{possession proof}
    $\pi = \mathsf{Sign}(\mathit{sk}_S, C_{\mathsf{reg}})$ where
    $C_{\mathsf{reg}}$ is a canonical byte sequence committing to
    $(\tau, S, T, K_S, t_{\mathsf{iat}}, \mathsf{nonce})$.
    \item $S$ submits $(\tau, S, T, K_S, \pi)$ to the registration
    endpoint.
    \item The server atomically (a) checks $\tau$ is unredeemed and
    unexpired, (b) marks $\tau$ redeemed, (c) verifies $\pi$ under
    $K_S$, and (d) writes a \texttt{KeyRecord} binding $K_S$ to
    $S$ under $T$ in state \texttt{Active}.
\end{enumerate}

\derived{} Step 2's possession proof prevents an adversary who
intercepts $\tau$ from substituting its own public key: a forged
$(\tau, S', T, K_{\mathsf{atk}}, \pi')$ still requires
$\mathit{sk}_{\mathsf{atk}}$ to produce $\pi'$, but binds the
registration to $K_{\mathsf{atk}}$ rather than the legitimate $K_S$.
The legitimate $S$, observing its token redeemed by another key, can
raise an alert and request a fresh token.

\subsection{The bootstrap invariant}

\begin{invariant}
\textbf{(Bootstrap separation.)}
\begin{itemize}[leftmargin=1.4em,itemsep=0.1em]
    \item A bootstrap token shall \emph{never} appear on any wire
    carrying an operation other than identity registration, and
    \emph{never} appear after the registration transaction has
    completed, whether successfully or not.
    \item A durable identity keypair's private component shall
    \emph{never} be generated, stored, or transmitted by any party
    other than the principal that owns the identity. Servers,
    intermediaries, backups, and logs shall not contain it.
\end{itemize}
These two prohibitions are symmetric. Violating either collapses the
separation.
\end{invariant}

\derived{} Violating the first prohibition reduces identity to a
bearer token, defeating G1 and G3. Violating the second turns the
violating party into a custodian indistinguishable from the
principal, defeating G3 (against the principal) and G6.

% ──────────────────────────────────────────────────────────────────────────
\section{Signing pipeline}
\label{sec:sign}

\subsection{Canonicalization}

\axiom{} For every invocation $I$, the canonical byte sequence $C(I)$
is computed once at the signing boundary (before any hop) and carried
unchanged through every intermediary until verification. Intermediaries
MUST NOT re-canonicalize; they MUST NOT re-order or re-encode any
field that contributes to $C(I)$.

\axiom{} $C(I)$ is domain-separated by a protocol version string so
that a signature valid under version $v$ is not accepted under a
different version.

\subsection{Signature production}

\derived{} To produce evidence for a binding rule $B$ over $I$:

\begin{itemize}[leftmargin=2em,itemsep=0.15em]
    \item The signer computes $\sigma = \mathsf{Sign}(\mathit{sk}, C(I))$
    using the private key corresponding to the identity the rule
    requires.
    \item The signer emits the tuple
    $(\sigma, K, \mathsf{kfp}(K), \mathsf{alg}(K))$ as companion metadata
    to $I$, where $\mathsf{kfp}(K)$ is the fingerprint and
    $\mathsf{alg}(K)$ is the algorithm identifier.
    \item Multiple binding rules contribute independent tuples; they
    are not combined, aggregated, or threshold-signed in v1.
\end{itemize}

\subsection{Signature suite}

\axiom{} The v1 signature suite is:

\begin{center}
\renewcommand{\arraystretch}{1.25}
\begin{tabularx}{\linewidth}{L{0.18\linewidth} L{0.22\linewidth} X}
\toprule
\textbf{Role} & \textbf{Algorithm} & \textbf{Reference} \\
\midrule
Primary & Ed25519 & RFC~8032~\cite{rfc8032},
\cite{bernstein2012high} \\
Alternate & ECDSA-P256 (SHA-256) & FIPS~186-5~\cite{fips186-5} \\
Hash & SHA-256 & FIPS~180-4 \\
RNG & OS CSPRNG (\texttt{getrandom}, \texttt{BCryptGenRandom}) &
NIST SP~800-90A \\
\bottomrule
\end{tabularx}
\end{center}

\deferred{} Post-quantum signature suites (FIPS~204
ML-DSA~\cite{fips204}, FIPS~205 SLH-DSA) are accommodated by the
\texttt{algorithm} field in \texttt{KeyRecord} but not required in v1.
See \S\ref{sec:gaps}.

% ──────────────────────────────────────────────────────────────────────────
\section{Verification pipeline}
\label{sec:verify}

\subsection{Mandatory verification}

\begin{invariant}
\textbf{(Mandatory verification, fail-closed.)} The runtime shall
evaluate every incoming invocation against the verification procedure
of \S\ref{sec:verify}.2. An invocation for which the procedure does
not return \texttt{Accept} shall be rejected. There shall be no
configuration under which unverified invocations are silently
accepted. Migration accommodations, if needed by an operator, shall
take the form of \emph{audit classes} emitted alongside rejection,
never in place of it.
\end{invariant}

\subsection{Verification procedure}

\axiom{} The runtime, for each incoming invocation $I$ with
accompanying signature tuples $\Sigma$, executes:

\begin{enumerate}[leftmargin=2em,itemsep=0.15em]
    \item \textbf{Canonical reconstruction.} Reconstruct $C(I)$ from
    the invocation. If any required field of $I$ is absent or
    unparseable, \texttt{Reject(MalformedCanonical)}.
    \item \textbf{Payload integrity.} Verify that
    $H(\mathsf{args})$ equals the hash committed to in $C(I)$.
    Else \texttt{Reject(PayloadIntegrity)} (G2).
    \item \textbf{Timestamp freshness.} Verify
    $|\mathsf{now} - t_{\mathsf{iat}}| \le \Delta_{\mathsf{skew}}$.
    Else \texttt{Reject(StaleOrFuture)} (G4).
    \item \textbf{Nonce reservation.} Atomically reserve
    $\mathsf{nonce}$ in the per-tenant nonce registry for a window of
    at least $2\Delta_{\mathsf{skew}}$. If already reserved,
    \texttt{Reject(Replay)} (G4).
    \item \textbf{Binding coverage.} Compute the set of binding rules
    required to authorize $\Phi(I)$ (\S\ref{sec:binding}). For each
    rule, verify the corresponding tuple in $\Sigma$: the signature
    verifies under the claimed public key; the public key resolves in
    the key registry to a record of state \texttt{Active}, subject
    matching the rule, tenant matching $T$, algorithm matching the
    algorithm indicator; and the expiry, if any, is not past.
    Any mismatch yields \texttt{Reject(BindingUnsatisfied)} (G1).
    \item \textbf{Non-substitutability.} Verify that no single
    signature is being applied to more than one binding rule
    (invariant of \S\ref{sec:binding}). Else
    \texttt{Reject(BindingCollision)}.
    \item \textbf{Audit commit.} Emit an audit record containing
    $\mathsf{req\_id}$, $T$, $P$, $N$, $A$, $f$, $\mathsf{kfp}(\cdot)$
    for each verifying key, and each signature $\sigma$. This record
    is the non-repudiation artifact (G3, G5).
    \item \texttt{Accept}. Dispatch $I$ to execution.
\end{enumerate}

\subsection{No opt-in mode}

\derived{} Verification is not gated on the presence of any
opt-in signal in the invocation (no ``if-metadata-present-then-verify''
shortcut). Invocations lacking signature tuples are not ``legacy
clients'' to the runtime; they are rejected. Operators preparing for
deployment are expected to upgrade clients before enabling the runtime,
not to leave enforcement conditional on the client's cooperation.

% ──────────────────────────────────────────────────────────────────────────
\section{Key lifecycle}
\label{sec:keys}

\subsection{Generation}

\axiom{} Private keys are generated on the party that owns them, using
the host operating system's cryptographic RNG:

\begin{itemize}[leftmargin=2em,itemsep=0.15em]
    \item Node keys: generated on the node at first boot, or at any
    subsequent rotation.
    \item Strong principal keys: generated in the principal's
    ultimate trust boundary (browser WebCrypto~\cite{webcrypto},
    CLI-local, hardware token).
    \item Weak principal keys (service accounts): generated on the
    hosting service at bootstrap, stored per the service's own
    secrets-management policy.
    \item Trust root keys: generated by the operator at runtime
    initialization; rotation is a privileged operator action.
\end{itemize}

\axiom{} Deterministic key derivation from public or low-entropy inputs
is prohibited. A key derived from inputs knowable to the adversary is
not a key.

\subsection{Storage}

\axiom{} Private keys at rest are protected by, minimally, host
operating-system access control (POSIX mode \texttt{0600} or the
equivalent ACL). Implementations SHOULD additionally support:
passphrase encryption (scrypt~\cite{rfc7914} key derivation, authenticated
symmetric encryption), OS key storage (macOS Keychain, Linux Secret
Service, Windows DPAPI), hardware-backed storage (TPM, Secure Enclave,
YubiKey, HSM).

\derived{} The signing interface is parameterized over a
\emph{signer} abstraction so that storage backends can be added
without changes to canonicalization or verification.

\subsection{Rotation}

\axiom{} Rotation of an identity's keypair proceeds as follows:

\begin{enumerate}[leftmargin=2em,itemsep=0.15em]
    \item The identity generates a fresh keypair
    $(\mathit{sk}',\,K')$ per \S\ref{sec:keys}.1.
    \item The identity produces a \emph{rotation proof}
    $\pi_{\mathsf{rot}} = \mathsf{Sign}(\mathit{sk},
    C_{\mathsf{rot}})$ under the current active key $\mathit{sk}$,
    where $C_{\mathsf{rot}}$ canonically commits to $(S, T, K', K,
    t_{\mathsf{iat}}, \mathsf{nonce})$.
    \item The identity submits $(S, T, K', \pi_{\mathsf{rot}})$ to
    the runtime's identity service.
    \item The runtime verifies $\pi_{\mathsf{rot}}$ under $K$,
    atomically transitions the prior \texttt{KeyRecord} to state
    \texttt{Rotated}, inserts a new record with state \texttt{Active}
    and version incremented, and records the rotation in audit.
\end{enumerate}

\derived{} The rotation proof binds the new key to the prior key's
authority: only the current legitimate holder can authorize rotation.
An adversary who obtains $\mathit{sk}$ can rotate, but this is
detectable (the legitimate holder observes a rotation it did not
initiate) and recoverable (revoke both; re-bootstrap).

\subsection{Revocation}

\axiom{} Revocation transitions a \texttt{KeyRecord} from any state to
\texttt{Revoked}. Signatures produced under a revoked key are
rejected at step 5 of the verification procedure from the moment
revocation is visible to the runtime.

\derived{} Revocation is authorized by one of: (i) the identity
itself, via a signature by the current key (self-revocation);
(ii) the trust root, via a signature by $K_R$ (administrative
revocation); (iii) out-of-band operator action (break-glass), which
MUST emit a distinguished audit class.

\derived{} Revocation is eventually consistent across federated
runtimes; cross-runtime propagation is \deferred (\S\ref{sec:gaps}).

\subsection{Expiry}

\derived{} A \texttt{KeyRecord} with
$\mathtt{expires\_at} > 0$ and $\mathtt{expires\_at} \le \mathsf{now}$
is treated as if revoked at verification time, without requiring an
explicit revocation operation. Expiry policy (default per-tenant
lifetimes, renewal grace periods) is \deferred.

% ──────────────────────────────────────────────────────────────────────────
\section{Trust roots and host verification}
\label{sec:trust-root}

\axiom{} The runtime exposes its own public identity $K_R$ through a
trust-root service. Clients pin $K_R$ at first contact, analogous to
SSH's \texttt{known\_hosts} (RFC~4251~\cite{rfc4251} \S4.1) or
WebPKI certificate pinning (RFC~7469~\cite{rfc7469}).

\derived{} Pinning mitigates the Dolev--Yao adversary's ability to
impersonate the runtime: after first contact, a client rejects a
runtime presenting a different $K_R$ unless a trust-root rotation proof
is supplied.

\axiom{} Trust-root rotation is itself an operation subject to this
specification: a rotation proof signed by the outgoing $K_R$ over the
incoming $K_R'$, countersigned as operator policy requires.

\deferred{} Multi-root endorsement (``the operator delegates trust to
a certificate authority''), cross-runtime trust in federation, and
the ceremony for the very first $K_R$ (bootstrap of the root) are
out of v1 scope.

% ──────────────────────────────────────────────────────────────────────────
\section{Protocol surface}
\label{sec:protocol}

This section states the minimum protocol surface the specification
requires. It is intentionally algorithm-agile and wire-format-agnostic
beyond the canonicalization requirements of \S\ref{sec:sign}.

\subsection{Invocation envelope}

\axiom{} Every invocation carries an \emph{envelope} with, minimally,
the fields of tuple $I$ (\S\ref{sec:statement}.1): tenant, principals,
node, resource, function, arguments, timestamp, request id, nonce.
Absence of any mandatory field at verification time yields
\texttt{Reject(MalformedCanonical)}.

\subsection{Signature metadata}

\axiom{} Signature tuples travel alongside the envelope, in metadata
clearly separable from the envelope (so the envelope's canonical form
is not perturbed). For each binding rule $B$ claimed, metadata
includes: the signature $\sigma$, the public key $K$, the key
fingerprint $\mathsf{kfp}(K)$, the algorithm identifier
$\mathsf{alg}(K)$, and the binding rule identifier
$B \in \{B1, B2, B3, \ldots\}$.

\subsection{Identity-service RPCs}

\axiom{} The runtime exposes at least the following identity-service
operations:

\begin{itemize}[leftmargin=2em,itemsep=0.15em]
    \item \texttt{RegisterKey}: the bootstrap registration of
    \S\ref{sec:bootstrap}.3.
    \item \texttt{RotateKey}: the rotation of \S\ref{sec:keys}.3.
    \item \texttt{RevokeKey}: the revocation of \S\ref{sec:keys}.4.
    \item \texttt{GetTrustRoot}: the public identity of the runtime.
    \item \texttt{RotateTrustRoot}: administrative trust-root rotation.
    \item \texttt{VerifySignature}: a standalone verification endpoint
    for offline audit verification.
\end{itemize}

\axiom{} Every identity-service RPC is itself an invocation subject to
this specification. Bootstrap registration is the single base case:
it is authorized by the bootstrap token rather than a signature,
exactly because it creates the signing capability that all subsequent
RPCs will require.

% ──────────────────────────────────────────────────────────────────────────
\section{Audit}
\label{sec:audit}

\axiom{} Every \texttt{Accept} outcome of the verification procedure
produces an audit record containing: $\mathsf{req\_id}$, tenant,
principal set with per-principal $\mathsf{kfp}$ and B2 mode
(strong / weak), node with its $\mathsf{kfp}$, resource and function,
argument hash, timestamp, binding rules satisfied, and each signature
$\sigma$.

\axiom{} Every \texttt{Reject} outcome produces an audit record
containing the same identifying fields (to the extent they parsed
successfully) plus the \texttt{Reject} reason code.

\derived{} The record's inclusion of raw signatures and public-key
fingerprints is what establishes G3 (non-repudiation) and G5 (forward
security): any later verifier with access to the historical key
registry can independently re-verify accepted records.

\axiom{} Audit records are append-only; amendments take the form of
new records, not mutations of old ones. Retention policy, storage
location, and access control are out of v1 scope.

% ──────────────────────────────────────────────────────────────────────────
\section{Construction of security properties}
\label{sec:construction}

\derived{} This section traces how G1--G6 follow from the mechanisms
above. The traces are informal but machine-checkable against the
preceding sections.

\begin{center}
\renewcommand{\arraystretch}{1.3}
\begin{xltabular}{\linewidth}{l X}
\toprule
\textbf{Goal} & \textbf{Construction} \\
\midrule
\endfirsthead
\toprule
\textbf{Goal} & \textbf{Construction} \\
\midrule
\endhead

G1 & Step 5 of \S\ref{sec:verify}.2 rejects invocations whose
identity claims lack a signature that verifies under a registered
\texttt{Active} key. Under \S\ref{sec:adversary}'s Ed25519 assumption,
only the private-key holder can produce such a signature.
\S\ref{sec:keys}.1 requires private keys be generated by the holder
and forbids derivation from public inputs; \S\ref{sec:bootstrap}.4
forbids custody transfer. \\

G2 & Step 2 recomputes $H(\mathsf{args})$ from the on-wire payload
and requires equality with the value inside $C(I)$. Canonicalization
covers every other routing-relevant field (\S\ref{sec:statement}.1).
Modification by a third party yields an invalid signature under G1's
construction. \\

G3 & Step 7 commits the signatures, public-key fingerprints, and
invocation fields to audit. A later verifier, with access to the
public-key registry at the time of acceptance, can re-run steps
5--6 offline against the audited record. Strong vs. weak B2 is
recorded to make the scope of the non-repudiation claim explicit. \\

G4 & Steps 3 and 4 combine: within one skew window, the nonce
registry rejects replays; outside it, timestamp freshness rejects
captured-and-resubmitted invocations. The window is per tenant to
prevent cross-tenant nonce collision-space exhaustion. \\

G5 & Audit records (\S\ref{sec:audit}) include the key fingerprint at
acceptance time, not a dynamic reference to the key record. Later
revocation updates the registry (rejecting new invocations) but leaves
historical records semantically unchanged. Revocation under
\S\ref{sec:keys}.4 records its own audit entry, so the transition
$t'_{\mathsf{revoked}}$ is itself non-repudiable. \\

G6 & \S\ref{sec:binding}'s invariant forbids binding substitution.
\S\ref{sec:identity}.1's distinction between strong and weak
principal identity prevents a weak signature from silently standing
in for a strong one: the audit record surfaces the mode, and
authorization policy above this layer may require strong B2 for
sensitive operations. Delegation (B4) is explicitly \deferred, not
tacitly absorbed into B2. \\

\bottomrule
\end{xltabular}
\end{center}

% ──────────────────────────────────────────────────────────────────────────
\section{Open extensions}
\label{sec:gaps}

These extensions are compatible with the v1 design and are
deliberately deferred.

\begin{enumerate}[leftmargin=2em,itemsep=0.3em]
    \item \textbf{Hardware- and OS-backed key storage.} Secure Enclave,
    TPM, YubiKey, HSM, macOS Keychain, Windows DPAPI, Linux Secret
    Service. Accommodated by the signer abstraction of
    \S\ref{sec:keys}.2; no protocol change.

    \item \textbf{Post-quantum signatures.} NIST FIPS~204 ML-DSA and
    FIPS~205 SLH-DSA~\cite{fips204}. Accommodated by
    \texttt{algorithm} in \texttt{KeyRecord}; a PQC-only deployment or
    hybrid deployment (classical + PQC) is additive.

    \item \textbf{Binding rule B4 (delegation).} A principal $A$
    authorizing a principal $B$ to act on its behalf in a bounded
    scope and for a bounded time. Candidate designs include
    macaroons~\cite{birgisson2014macaroons}, capability
    URLs~\cite{capurl}, biscuits, and SPIFFE federation tokens. v1
    explicitly forbids delegation to avoid committing to a design
    prematurely.

    \item \textbf{Per-tenant policy.} Skew window, nonce TTL,
    required-mode for principals (strong-B2-only), accepted algorithms,
    key expiry policy. v1 uses runtime-global defaults.

    \item \textbf{Cross-runtime federation identity.} Revocation
    propagation, cross-tenant naming, federation-level trust roots.
    SPIFFE~\cite{spiffe} trust-bundle federation is the reference
    design.

    \item \textbf{Workload identity.} Non-interactive pairing for
    automation (CI runners, sidecars, initContainers) along the lines
    of SPIFFE \texttt{WorkloadAPI} and Kubernetes projected service
    account tokens.

    \item \textbf{Anonymous credentials and unlinkability.}
    Group-signature-based and attribute-based constructions for cases
    where the principal wishes not to be linked across invocations.
    Orthogonal to v1's non-repudiation goal, hence \deferred.

    \item \textbf{Threshold signatures.} A binding rule that requires
    $k$-of-$n$ key holders to jointly sign. Useful for high-stakes
    operations (trust-root rotation, cross-tenant delegation). Out of
    v1 because the additional protocol complexity is not justified by
    v1 use cases.

    \item \textbf{Formal verification.} The specification is amenable
    to mechanized proof in Tamarin~\cite{tamarin} or ProVerif. A
    formal model is out of v1 but the canonicalization and binding
    rules are structured to make formalization tractable.
\end{enumerate}

% ──────────────────────────────────────────────────────────────────────────
\section{Conclusion}

The load-bearing content of this specification is three items.

\emph{Mandatory verification, fail-closed} (\S\ref{sec:verify}):
every invocation is evaluated against a fixed procedure; invocations
that do not satisfy it are rejected; there is no opt-in shortcut.

\emph{Binding semantics} (\S\ref{sec:binding}): node-origin and
principal-authorization are distinct claims, each requiring its own
signature; neither substitutes for the other; conjunction is the only
way to claim both.

\emph{Bootstrap separation} (\S\ref{sec:bootstrap}): single-use
authorization tokens and durable identity keypairs play non-
interchangeable roles; the private component of a durable keypair
lives only on the party that owns the identity.

Everything else---the signing pipeline, the verification procedure,
the key lifecycle, the audit record---follows mechanically from these
three. An implementation that satisfies all three satisfies
G1 through G6 under the adversary model of \S\ref{sec:adversary}.
An implementation that relaxes any of them does not.

Extensions enumerated in \S\ref{sec:gaps} are additive: they augment
this specification without revisiting its axioms. The specification
is designed to be conformant-once, extended-many-times.

% ──────────────────────────────────────────────────────────────────────────
% References belong with the normative body, BEFORE \appendix, so they
% are not visually grouped with the conformance appendix. `thebibliography`
% in article.cls emits its own \section*{\refname} heading, so we do NOT
% add a manual \section*{References} above it (that would print the title
% twice). We only register a TOC entry.
\addcontentsline{toc}{section}{References}
\markboth{References}{References}

\begin{thebibliography}{99}

\bibitem{dolev1983security}
D.~Dolev and A.~Yao.
\newblock On the security of public key protocols.
\newblock \emph{IEEE Transactions on Information Theory}, 29(2):198--208,
1983.

\bibitem{rfc4251}
T.~Ylonen and C.~Lonvick (Eds.).
\newblock The Secure Shell (SSH) Protocol Architecture.
\newblock RFC 4251, IETF, 2006.

\bibitem{rfc4252}
T.~Ylonen and C.~Lonvick (Eds.).
\newblock The Secure Shell (SSH) Authentication Protocol.
\newblock RFC 4252, IETF, 2006.

\bibitem{rfc7469}
C.~Evans, C.~Palmer, and R.~Sleevi.
\newblock Public Key Pinning Extension for HTTP.
\newblock RFC 7469, IETF, 2015.

\bibitem{rfc7914}
C.~Percival and S.~Josefsson.
\newblock The scrypt Password-Based Key Derivation Function.
\newblock RFC 7914, IETF, 2016.

\bibitem{rfc8032}
S.~Josefsson and I.~Liusvaara.
\newblock Edwards-Curve Digital Signature Algorithm (EdDSA).
\newblock RFC 8032, IETF, 2017.

\bibitem{rfc8446}
E.~Rescorla.
\newblock The Transport Layer Security (TLS) Protocol Version 1.3.
\newblock RFC 8446, IETF, 2018.

\bibitem{rfc8785}
A.~Rundgren, B.~Jordan, and S.~Erdtman.
\newblock JSON Canonicalization Scheme (JCS).
\newblock RFC 8785, IETF, 2020.

\bibitem{bernstein2012high}
D.~J. Bernstein, N.~Duif, T.~Lange, P.~Schwabe, and B.-Y. Yang.
\newblock High-speed high-security signatures.
\newblock \emph{Journal of Cryptographic Engineering}, 2(2):77--89, 2012.

\bibitem{fips186-5}
National Institute of Standards and Technology.
\newblock Digital Signature Standard (DSS).
\newblock FIPS PUB 186-5, 2023.

\bibitem{fips204}
National Institute of Standards and Technology.
\newblock Module-Lattice-Based Digital Signature Standard.
\newblock FIPS PUB 204 (ML-DSA), 2024.

\bibitem{spiffe}
SPIFFE Authors.
\newblock The SPIFFE Identity and Verifiable Identity Document (SPIFFE
ID, SVID).
\newblock SPIFFE Specifications v1.0+, \url{https://spiffe.io/}.

\bibitem{webauthn}
D.~Balfanz et~al.
\newblock Web Authentication: An API for accessing Public Key
Credentials, Level~2.
\newblock W3C Recommendation, 2021.

\bibitem{webcrypto}
M.~Watson (Ed.).
\newblock Web Cryptography API.
\newblock W3C Recommendation, 2017.

\bibitem{birgisson2014macaroons}
A.~Birgisson, J.~G. Politz, {\'U}.~Erlingsson, A.~Taly, M.~Vrable, and
M.~Lentczner.
\newblock Macaroons: Cookies with contextual caveats for decentralized
authorization in the cloud.
\newblock In \emph{NDSS}, 2014.

\bibitem{capurl}
M.~S. Miller.
\newblock Robust composition: Towards a unified approach to access
control and concurrency control.
\newblock PhD thesis, Johns Hopkins University, 2006.

\bibitem{tamarin}
S.~Meier, B.~Schmidt, C.~Cremers, and D.~Basin.
\newblock The Tamarin Prover for the symbolic analysis of security
protocols.
\newblock In \emph{CAV}, 2013.

\end{thebibliography}

% ──────────────────────────────────────────────────────────────────────────
\appendix
\section{Conformance of the present codebase}
\label{app:conformance}

This appendix is operational, not normative. It identifies the
deviations of the current EasyNet codebase (as observed at the time
of this edition) from the specification above. Each deviation is
tagged \defect, cited to file path and line number, and mapped to the
section of the specification it violates.

\begin{notebox}
The body of this document would remain valid, as a reference
specification, if this appendix were deleted. The appendix exists to
direct implementation work, not to influence the design.
\end{notebox}

\subsection{Summary}

\defect{} The present runtime accepts invocations that lack any
identity evidence, by virtue of an opt-in branch in the verifier
(violates \S\ref{sec:verify}). The three identity layers
(\S\ref{sec:identity}) exist as separate artifacts---client-side
bearer token, backend-side JWT, runtime-side keypair infrastructure---
none of which is transported end-to-end to the verification boundary.
The runtime executes in a state operationally indistinguishable from
``unauthenticated'' while appearing, to a surface audit, to be
authenticated. The specific loci are enumerated below.

\subsection{Non-conformance catalog}

Each entry names the deviation, the file locations at which it
manifests, and the section of the specification it violates. File
paths are given as verbatim identifiers to be resolvable without
ambiguity; they are not links.

\paragraph{C1 \textemdash{} opt-in verification}
\textbf{Violates:} \S\ref{sec:verify}.
The verifier returns \texttt{Ok(None)} when the expected
\texttt{easynet.*} metadata is absent, short-circuiting the
verification procedure and allowing the invocation to dispatch.
\\*
\emph{Location:} \texttt{core/runtime-rs/src/runtime/security.rs},
in \texttt{verify\_easynet\_invocation\_metadata} (the early return
when \texttt{easynet.resource\_uri} is absent).

\paragraph{C2 \textemdash{} tenant-only gate}
\textbf{Violates:} \S\ref{sec:verify}.
The sole authentication gate invoked from invoke handlers checks
envelope tenant presence only; no binding evaluation is performed.
\\*
\emph{Location:} \texttt{core/runtime-rs/src/runtime/mod.rs},
in \texttt{require\_envelope\_tenant}.

\paragraph{C3 \textemdash{} backend interface has no principal slot}
\textbf{Violates:} \S\ref{sec:identity}.
The backend's \texttt{CallMcpTool} (and
\texttt{CallMcpToolStream}) interface has no parameter for a
principal identifier; user identity captured from the JWT cannot
reach the runtime.
\\*
\emph{Location:} \texttt{backend/internal/axon/client.go},
interface definition.

\paragraph{C4 \textemdash{} principal dropped at Axon call site}
\textbf{Violates:} \S\ref{sec:identity}.
The backend ability-invoke logic captures \texttt{uid} from the
request context, then calls the Axon client without transporting it.
\\*
\emph{Location:}
\texttt{backend/internal/logic/ability/invokeAbilityLogic.go}.

\paragraph{C5 \textemdash{} \texttt{credential\_token} conflates roles}
\textbf{Violates:} \S\ref{sec:bootstrap}.
A single field \texttt{credential\_token} is used both as bootstrap
proof (issued and redeemed at pairing) and as durable identity
(stored at rest and replayed at every startup), collapsing the
distinction the specification requires.
\\*
\emph{Locations:} \texttt{EasyNet-Cli/src/shared/config.rs}
(the \texttt{credential\_token} field of \texttt{Credentials});
\texttt{backend/internal/logic/device/validatePairingLogic.go}
(issuance of plaintext token);
\texttt{EasyNet-Cli/src/cli/start.rs} (replay on every startup).

\paragraph{C6 \textemdash{} deterministic-derivation ``keys''}
\textbf{Violates:} \S\ref{sec:keys}.
The client SDK derives an Ed25519 seed deterministically from
\texttt{tenant\_id} and \texttt{subject\_id} using a public salt,
producing a private key recoverable by any party with access to the
SDK source.
\\*
\emph{Location:}
\texttt{core/runtime-rs/client-sdk/src/domain/easynet/semantic.rs},
in \texttt{derive\_subject\_auth}.

\paragraph{C7 \textemdash{} envelopes omit \texttt{easynet} metadata}
\textbf{Violates:} \S\S\ref{sec:sign}, \ref{sec:verify}.
CLI client envelopes sent to the runtime have the \texttt{easynet}
metadata field unset; no signature tuple is produced at the signing
boundary.
\\*
\emph{Location:}
\texttt{core/runtime-rs/client-sdk/src/infra/client.rs}, envelope
construction in the \texttt{heartbeat} and \texttt{invoke} paths,
where the \texttt{easynet} field is left \texttt{None}.

\paragraph{C8 \textemdash{} CLI call paths never sign}
\textbf{Violates:} \S\S\ref{sec:sign}, \ref{sec:verify}.
Heartbeat and ability-invoke call paths in the CLI do not construct
or attach signature metadata.
\\*
\emph{Locations:} \texttt{EasyNet-Cli/src/cli/heartbeat.rs};
\texttt{EasyNet-Cli/src/cli/invoke.rs}.

\subsection{Resolution phasing}

\derived{} The observed deviations can be resolved without a
specification change. Suggested phasing:

\begin{itemize}[leftmargin=2em,itemsep=0.2em]
    \item \textbf{Phase~R0 (code parity).} Resolve C6, C7, C8 on the
    signing side and C3, C4 on the transport side. The runtime still
    accepts unsigned invocations; audit distinguishes signed from
    unsigned classes. No invocation is rejected in this phase; the
    goal is to make every real caller capable of producing signatures.
    \item \textbf{Phase~R1 (enforcement).} Resolve C1, C2. The
    verifier becomes mandatory; unsigned invocations are rejected.
    Audit data from R0 is reviewed to confirm the unsigned class is
    empty before this phase begins.
    \item \textbf{Phase~R2 (cleanup).} Resolve C5 by removing
    \texttt{credential\_token} from persisted credentials and from
    the backend schema; the device's durable identity is the keypair
    alone.
\end{itemize}

\derived{} Phasing is an operational choice; the specification itself
requires only the end state (all of C1--C8 resolved). Permanent
retention of the R0 state is not a valid long-term configuration.

\vfill
\hrule
\vspace{0.4em}
\noindent\small
\textit{The body of this document (\S\S\ref{sec:reading-guide}--\ref{sec:gaps})
is a reference specification, design-first and implementation-agnostic.
The appendices analyze a specific codebase's conformance to it. The two
concerns are kept separate so the specification remains useful to future
clean-room implementations without inheriting the contingencies of the
present codebase.}

\end{document}
